\textbf{Proof.}
Every \(s\in\mathcal Z_n\) is balanced within each quadruple, so
\[
s\in\mathcal H
=\bigoplus_{\ell=1}^q\mathbf1_4^\perp,
\qquad \dim(\mathcal H)=3q=d. \tag{1}
\]
Also every coordinate of \(s\) is \(\pm1\), whence
\[
\operatorname{tr}(ss^\top)=\|s\|_2^2=n. \tag{2}
\]
Let \(\operatorname{Sym}(\mathcal H)\) denote the symmetric operators on \(\mathcal H\).  It has dimension
\(D=d(d+1)/2\).  The matrices \(ss^\top\) all lie in the affine hyperplane of \(\operatorname{Sym}(\mathcal H)\) having trace \(n\), whose dimension is at most \(D-1\).

Write an optimal second moment as
\[
Q^\star=\sum_{s\in\mathscr S}\alpha_s^star ss^\top,
\qquad \alpha^star\in\Delta_m, \tag{3}
\]
using the mass representation in \(\mathrm{thm:higher-rank-menu-exact-program}\).  We give the support reduction explicitly.  If more than \(D\) coefficients in (3) are positive, the corresponding points in an affine space of dimension at most \(D-1\) are affinely dependent.  Hence there are numbers \(c_s\), not all zero and supported on the positive coefficients, such that
\[
\sum_sc_s=0,
\qquad
\sum_sc_sss^\top=0. \tag{4}
\]
After changing the common sign if necessary, some \(c_s>0\).  Set
\[
\varepsilon=\min_{c_s>0}\frac{\alpha_s^star}{c_s}>0,
\qquad
\alpha_s'=\alpha_s^star-\varepsilon c_s. \tag{5}
\]
Equations (4)--(5) preserve total mass and the matrix in (3), keep every coefficient nonnegative, and make at least one previously positive coefficient zero.  Repeating finitely many times yields \(\widetilde\alpha\in\Delta_m\) with
\[
Q^\star=\sum_s\widetilde\alpha_sss^\top,
\qquad
|\operatorname{supp}(\widetilde\alpha)|\le\min\{m,D\}. \tag{6}
\]

Define the blinded law
\[
\widetilde\pi(s)=\widetilde\pi(-s)=\frac12\widetilde\alpha_s,
\qquad s\in\mathscr S. \tag{7}
\]
It belongs to \(\Pi_n\), and (6)--(7) give \(Q_{\widetilde\pi}=Q^\star\).  Substituting \(m=6^q/2\), \(d=3q\), and \(D=d(d+1)/2\) gives the claimed sign-pair bound
\[
\min\left\{\frac{6^q}{2},\frac{3q(3q+1)}2\right\}. \tag{8}
\]
Splitting each positive sign-pair mass uses at most twice as many assignments.  Every generalized risk and every constraint of \((\mathrm P)\) depends on the law only through \(Q_\pi\), so \(\widetilde\pi\) has exactly the same risks and commitment value as the original optimizer.  The reduction starts from \(Q^\star\); it does not furnish that matrix or locate its sparse decomposition without access to the exponentially enumerated columns.  Hence it is a support-existence result, not a polynomial-time algorithm. \(\square\)